\subsubsection{Infrastructure Adapters}

\paragraph{PostgresAdapter}
\vspace{\baselineskip}\noindent\textbf{Descrizione della classe:} implementa le porte IUserRepository, ISessionManager, IOrderRepository e ITicketRepository, interagendo direttamente con PostgreSQL tramite psycopg2 e connection pooling. Essendo un Adapter, svolge il ruolo cruciale di trasformare i tuple-record grezzi del database in veri e propri oggetti applicativi e DTO richiesti dalle interfacce del dominio.

\vspace{\baselineskip}\noindent\textbf{Descrizione degli attributi:}
\begin{itemize}
    \item \texttt{pool}: istanza di \texttt{SimpleConnectionPool} per la gestione ottimizzata delle connessioni al database.
\end{itemize}

\vspace{\baselineskip}\noindent\textbf{Descrizione dei metodi:}
\begin{itemize}
    \item \texttt{find\_by\_email(email)}: implementa IUserRepository — recupera utente per email.
    \item \texttt{get\_active\_session(user\_id)}: implementa ISessionManager — recupera sessione attiva.
    \item \texttt{create\_order(cod\_cli, user\_id, session\_id, items)}: implementa IOrderRepository — crea ordine.
    \item \texttt{create\_ticket(session\_id, cod\_cli)}: implementa ITicketRepository — crea ticket.
    \item \texttt{search\_products(query, cod\_cli, limit)}: ricerca prodotti nell'assortimento cliente.
    \item \texttt{find\_product\_by\_name(name, cod\_cli)}: ricerca prodotto per nome testuale.
\end{itemize}

\begin{figure}[H]
    \centering
    \includegraphics[width=0.52\textwidth]{uml_backend/35_postgresadapter.png}
    \caption{UML della classe \texttt{PostgresAdapter}.}
\end{figure}

\paragraph{OpenAIAdapter}
\vspace{\baselineskip}\noindent\textbf{Descrizione della classe:} realizza il contratto \texttt{IAIClient} mappando chiamate di dominio verso le API OpenAI (GPT-5.4, Whisper, Embeddings, Vision).

\vspace{\baselineskip}\noindent\textbf{Descrizione degli attributi:} il diagramma non esplicita attributi di configurazione (caricati da config).

\vspace{\baselineskip}\noindent\textbf{Descrizione dei metodi:}
\begin{itemize}
    \item \texttt{transcribe\_audio(audio\_file, filename)}: trascrizione audio via Whisper con prompt specifico per il dominio Ergon.
    \item \texttt{interpret\_intent(text, context, system\_prompt, model)}: invoca il modello LLM con sistema prompt e contesto.
    \item \texttt{generate\_embedding(text, model)}: genera un embedding vettoriale per il testo fornito.
    \item \texttt{analyze\_image(image\_base64, conversation\_history, system\_prompt)}: analizza immagini con GPT-5.4 Vision e restituisce OCR strutturato.
\end{itemize}

\begin{figure}[H]
    \centering
    \includegraphics[width=0.62\textwidth]{uml_backend/36_openaiadapter.png}
    \caption{UML della classe \texttt{OpenAIAdapter}.}
\end{figure}
